\documentclass{amsart}
% For dealing with references we use the comment environment
\usepackage{verbatim}
\newenvironment{reference}{\comment}{\endcomment}
%\newenvironment{reference}{}{}
\newenvironment{slogan}{\comment}{\endcomment}
\newenvironment{history}{\comment}{\endcomment}

% For commutative diagrams we use Xy-pic
\usepackage[all]{xy}

% We use 2cell for 2-commutative diagrams.
\xyoption{2cell}
\UseAllTwocells

% We use multicol for the list of chapters between chapters
\usepackage{multicol}

% This is generally recommended for better output
\usepackage{lmodern}
\usepackage[T1]{fontenc}

% For cross-file-references

% Package for hypertext links:
\usepackage{hyperref}

% For any local file, say "hello.tex" you want to link to please

% Theorem environments.
%
\theoremstyle{plain}
\newtheorem{theorem}[subsection]{Theorem}
\newtheorem{proposition}[subsection]{Proposition}
\newtheorem{lemma}[subsection]{Lemma}

\theoremstyle{definition}
\newtheorem{definition}[subsection]{Definition}
\newtheorem{example}[subsection]{Example}
\newtheorem{exercise}[subsection]{Exercise}
\newtheorem{situation}[subsection]{Situation}

\theoremstyle{remark}
\newtheorem{remark}[subsection]{Remark}
\newtheorem{remarks}[subsection]{Remarks}

\numberwithin{equation}{subsection}

% Macros
%
\def\lim{\mathop{\mathrm{lim}}\nolimits}
\def\colim{\mathop{\mathrm{colim}}\nolimits}
\def\Spec{\mathop{\mathrm{Spec}}}
\def\Hom{\mathop{\mathrm{Hom}}\nolimits}
\def\Ext{\mathop{\mathrm{Ext}}\nolimits}
\def\SheafHom{\mathop{\mathcal{H}\!\mathit{om}}\nolimits}
\def\SheafExt{\mathop{\mathcal{E}\!\mathit{xt}}\nolimits}
\def\Sch{\mathit{Sch}}
\def\Mor{\mathop{\mathrm{Mor}}\nolimits}
\def\Ob{\mathop{\mathrm{Ob}}\nolimits}
\def\Sh{\mathop{\mathit{Sh}}\nolimits}
\def\NL{\mathop{N\!L}\nolimits}
\def\CH{\mathop{\mathrm{CH}}\nolimits}
\def\proetale{{pro\text{-}\acute{e}tale}}
\def\etale{{\acute{e}tale}}
\def\QCoh{\mathit{QCoh}}
\def\Ker{\mathop{\mathrm{Ker}}}
\def\Im{\mathop{\mathrm{Im}}}
\def\Coker{\mathop{\mathrm{Coker}}}
\def\Coim{\mathop{\mathrm{Coim}}}

% Boxtimes
%
\DeclareMathSymbol{\boxtimes}{\mathbin}{AMSa}{"02}

%
% Macros for moduli stacks/spaces
%
\def\QCohstack{\mathcal{QC}\!\mathit{oh}}
\def\Cohstack{\mathcal{C}\!\mathit{oh}}
\def\Spacesstack{\mathcal{S}\!\mathit{paces}}
\def\Quotfunctor{\mathrm{Quot}}
\def\Hilbfunctor{\mathrm{Hilb}}
\def\Curvesstack{\mathcal{C}\!\mathit{urves}}
\def\Polarizedstack{\mathcal{P}\!\mathit{olarized}}
\def\Complexesstack{\mathcal{C}\!\mathit{omplexes}}
% \Pic is the operator that assigns to X its picard group, usage \Pic(X)
% \Picardstack_{X/B} denotes the Picard stack of X over B
% \Picardfunctor_{X/B} denotes the Picard functor of X over B
\def\Pic{\mathop{\mathrm{Pic}}\nolimits}
\def\Picardstack{\mathcal{P}\!\mathit{ic}}
\def\Picardfunctor{\mathrm{Pic}}
\def\Deformationcategory{\mathcal{D}\!\mathit{ef}}

\setlength{\emergencystretch}{3em}
\begin{document}
\title[Stacks Project, Tag 08SZ]{395 --- The transitivity triangle for cotangent complexes of sheaves of rings}
\maketitle
\noindent Copyright (C) 2005--2025 Johan de Jong. Stacks Project authors. Source: \href{https://stacks.math.columbia.edu/tag/08SZ}{Tag 08SZ}.

\noindent Editorial difficulty note (not author text). Difficulty H2: The proof constructs compatible standard resolutions and an exact simplicial differential sequence when the second map is injective. It identifies the final term through two quasi-isomorphisms, then uses sheafification and a product-ring replacement to handle arbitrary homomorphisms. This is the global construction of the canonical triangle, with nontrivial resolution compatibility..

\noindent Editorial provenance note (not author text). History: excerpt from revision \texttt{a0444\allowbreak 6e57e\allowbreak c1fbc\allowbreak 252a8\allowbreak 71afc\allowbreak ec775\allowbreak 2fb28\allowbreak 07b14}, cotangent.tex, lines 3277--3351. Reference presentation was adapted; original-author context and core support blocks were retained or added. The mathematical wording is unchanged. Licensed under GNU FDL 1.2 or later; no invariant sections or cover texts. See ../licenses/COPYING. Context blocks reproduce author definitions and supporting results; they are not additional counted problems.

\begin{definition}
\label{definition-standard-resolution-sheaves-rings}
Let $\mathcal{C}$ be a site.
Let $\mathcal{A} \to \mathcal{B}$ be a homomorphism of sheaves of rings
on $\mathcal{C}$. The {\it standard resolution of $\mathcal{B}$ over
$\mathcal{A}$} is the augmentation
$\epsilon : \mathcal{P}_\bullet \to \mathcal{B}$
with terms
$$
\mathcal{P}_0 = \mathcal{A}[\mathcal{B}],\quad
\mathcal{P}_1 = \mathcal{A}[\mathcal{A}[\mathcal{B}]],\quad \ldots
$$
and maps as constructed above.
\end{definition}

\begin{definition}
\label{definition-cotangent-complex-morphism-sheaves-rings}
Let $\mathcal{C}$ be a site.
Let $\mathcal{A} \to \mathcal{B}$ be a homomorphism of sheaves of rings
on $\mathcal{C}$.
The {\it cotangent complex} $L_{\mathcal{B}/\mathcal{A}}$
is the complex of $\mathcal{B}$-modules associated to the
simplicial module
$$
\Omega_{\mathcal{P}_\bullet/\mathcal{A}}
\otimes_{\mathcal{P}_\bullet, \epsilon} \mathcal{B}
$$
where $\epsilon : \mathcal{P}_\bullet \to \mathcal{B}$
is the standard resolution of $\mathcal{B}$ over
$\mathcal{A}$. We usually think of $L_{\mathcal{B}/\mathcal{A}}$
as an object of $D(\mathcal{B})$.
\end{definition}

\begin{remark}
\label{remark-triangle}
We sketch an alternative, perhaps simpler, proof of the existence of
the fundamental triangle.
Let $A \to B \to C$ be ring maps and assume that $B \to C$ is injective.
Let $P_\bullet \to B$ be the standard resolution of $B$ over $A$ and
let $Q_\bullet \to C$ be the standard resolution of $C$ over $B$.
Picture
$$
\xymatrix{
P_\bullet : &
A[A[A[B]]] \ar[d]
\ar@<2ex>[r]
\ar@<0ex>[r]
\ar@<-2ex>[r]
&
A[A[B]] \ar[d]
\ar@<1ex>[r]
\ar@<-1ex>[r]
\ar@<1ex>[l]
\ar@<-1ex>[l]
&
A[B] \ar[d] \ar@<0ex>[l] \ar[r] &
B \\
Q_\bullet : &
A[A[A[C]]]
\ar@<2ex>[r]
\ar@<0ex>[r]
\ar@<-2ex>[r]
&
A[A[C]]
\ar@<1ex>[r]
\ar@<-1ex>[r]
\ar@<1ex>[l]
\ar@<-1ex>[l]
&
A[C] \ar@<0ex>[l] \ar[r] &
C
}
$$
Observe that since $B \to C$ is injective, the ring $Q_n$ is a
polynomial algebra over $P_n$ for all $n$. Hence we obtain a cosimplicial
object in $\mathcal{C}_{C/B/A}$ (beware reversal arrows).
Now set $\overline{Q}_\bullet = Q_\bullet \otimes_{P_\bullet} B$.
The key to the proof of Proposition \href{https://stacks.math.columbia.edu/tag/08QX}{proposition triangle (Tag 08QX)}
is to show that $\overline{Q}_\bullet$ is a resolution of $C$ over $B$.
This follows from Cohomology on Sites, Lemma
\href{https://stacks.math.columbia.edu/tag/08RX}{lemma O homology qis (Tag 08RX)}
applied to $\mathcal{C} = \Delta$, $\mathcal{O} = P_\bullet$,
$\mathcal{O}' = B$, and $\mathcal{F} = Q_\bullet$ (this uses that $Q_n$
is flat over $P_n$; see Cohomology on Sites, Remark
\href{https://stacks.math.columbia.edu/tag/08QD}{remark simplicial modules (Tag 08QD)} to relate simplicial modules
to sheaves). The key fact implies that the distinguished triangle of
Proposition \href{https://stacks.math.columbia.edu/tag/08QX}{proposition triangle (Tag 08QX)}
is the distinguished triangle associated to the short exact sequence
of simplicial $C$-modules
$$
0 \to
\Omega_{P_\bullet/A} \otimes_{P_\bullet} C \to
\Omega_{Q_\bullet/A} \otimes_{Q_\bullet} C \to
\Omega_{\overline{Q}_\bullet/B} \otimes_{\overline{Q}_\bullet} C \to 0
$$
which is deduced from the short exact sequences
$0 \to \Omega_{P_n/A} \otimes_{P_n} Q_n \to \Omega_{Q_n/A} \to
\Omega_{Q_n/P_n} \to 0$ of
Algebra, Lemma \href{https://stacks.math.columbia.edu/tag/031K}{lemma ses formally smooth (Tag 031K)}.
Namely, by Remark \href{https://stacks.math.columbia.edu/tag/08QI}{remark resolution (Tag 08QI)} and the key fact the complex on the
right hand side represents $L_{C/B}$ in $D(C)$.

\medskip\noindent
If $B \to C$ is not injective, then we can use the above to get a
fundamental triangle for $A \to B \to B \times C$. Since
$L_{B \times C/B} \to L_{B/B} \oplus L_{C/B}$ and
$L_{B \times C/A} \to L_{B/A} \oplus L_{C/A}$
are quasi-isomorphism in $D(B \times C)$
(Lemma \href{https://stacks.math.columbia.edu/tag/08SC}{lemma cotangent complex product (Tag 08SC)})
this induces the desired distinguished triangle in $D(C)$
by tensoring with the flat ring map $B \times C \to C$.
\end{remark}

\begin{remark}
\label{remark-explicit-map}
Let $A \to B \to C$ be ring maps with $B \to C$ injective.
Recall the notation $P_\bullet$, $Q_\bullet$, $\overline{Q}_\bullet$ of
Remark \ref{remark-triangle}.
Let $R_\bullet$ be the standard resolution of $C$ over $B$.
In this remark we explain how to get the canonical identification
of $\Omega_{\overline{Q}_\bullet/B} \otimes_{\overline{Q}_\bullet} C$
with $L_{C/B} = \Omega_{R_\bullet/B} \otimes_{R_\bullet} C$.
Let $S_\bullet \to B$ be the standard resolution of $B$ over $B$.
Note that the functoriality map $S_\bullet \to R_\bullet$ identifies
$R_n$ as a polynomial algebra over $S_n$ because $B \to C$ is injective.
For example in degree $0$ we have the map $B[B] \to B[C]$, in degree
$1$ the map $B[B[B]] \to B[B[C]]$, and so on. Thus
$\overline{R}_\bullet = R_\bullet \otimes_{S_\bullet} B$
is a simplicial polynomial algebra
over $B$ as well and it follows (as in Remark \ref{remark-triangle}) from
Cohomology on Sites, Lemma
\href{https://stacks.math.columbia.edu/tag/08RX}{lemma O homology qis (Tag 08RX)}
that $\overline{R}_\bullet \to C$ is a resolution. Since we have
a commutative diagram
$$
\xymatrix{
Q_\bullet \ar[r] & R_\bullet \\
P_\bullet \ar[u] \ar[r] & S_\bullet \ar[u] \ar[r] & B
}
$$
we obtain a canonical map
$\overline{Q}_\bullet = Q_\bullet \otimes_{P_\bullet} B \to
\overline{R}_\bullet$. Thus the maps
$$
L_{C/B} = \Omega_{R_\bullet/B} \otimes_{R_\bullet} C
\longrightarrow
\Omega_{\overline{R}_\bullet/B} \otimes_{\overline{R}_\bullet} C
\longleftarrow
\Omega_{\overline{Q}_\bullet/B} \otimes_{\overline{Q}_\bullet} C
$$
are quasi-isomorphisms (Remark \href{https://stacks.math.columbia.edu/tag/08QI}{remark resolution (Tag 08QI)}) and composing
one with the inverse of the other gives the desired identification.
\end{remark}

\begin{lemma}
\label{lemma-triangle-sheaves-rings}
Let $\mathcal{D}$ be a site. Let $\mathcal{A} \to \mathcal{B} \to \mathcal{C}$
be homomorphisms of sheaves of rings on $\mathcal{D}$.
There is a canonical distinguished triangle
$$
L_{\mathcal{B}/\mathcal{A}} \otimes_\mathcal{B}^\mathbf{L} \mathcal{C}
\to L_{\mathcal{C}/\mathcal{A}} \to L_{\mathcal{C}/\mathcal{B}} \to
L_{\mathcal{B}/\mathcal{A}} \otimes_\mathcal{B}^\mathbf{L} \mathcal{C}[1]
$$
in $D(\mathcal{C})$.
\end{lemma}

\begin{proof}
We will use the method described in
Remarks \ref{remark-triangle} and \ref{remark-explicit-map}
to construct the triangle; we will freely use the results mentioned there.
As in those remarks we first construct the triangle in case
$\mathcal{B} \to \mathcal{C}$ is an injective map of sheaves of rings.
In this case we set
\begin{enumerate}
\item $\mathcal{P}_\bullet$ is the standard resolution of $\mathcal{B}$
over $\mathcal{A}$,
\item $\mathcal{Q}_\bullet$ is the standard resolution of $\mathcal{C}$
over $\mathcal{A}$,
\item $\mathcal{R}_\bullet$ is the standard resolution of $\mathcal{C}$
over $\mathcal{B}$,
\item $\mathcal{S}_\bullet$ is the standard resolution of $\mathcal{B}$
over $\mathcal{B}$,
\item $\overline{\mathcal{Q}}_\bullet =
\mathcal{Q}_\bullet \otimes_{\mathcal{P}_\bullet} \mathcal{B}$, and
\item $\overline{\mathcal{R}}_\bullet =
\mathcal{R}_\bullet \otimes_{\mathcal{S}_\bullet} \mathcal{B}$.
\end{enumerate}
The distinguished triangle is the distinguished triangle associated
to the short exact sequence
of simplicial $\mathcal{C}$-modules
$$
0 \to
\Omega_{\mathcal{P}_\bullet/\mathcal{A}}
\otimes_{\mathcal{P}_\bullet} \mathcal{C} \to
\Omega_{\mathcal{Q}_\bullet/\mathcal{A}}
\otimes_{\mathcal{Q}_\bullet} \mathcal{C} \to
\Omega_{\overline{\mathcal{Q}}_\bullet/\mathcal{B}}
\otimes_{\overline{\mathcal{Q}}_\bullet} \mathcal{C} \to 0
$$
The first two terms are equal to the first two terms of the triangle
of the statement of the lemma. The identification of the last term with
$L_{\mathcal{C}/\mathcal{B}}$ uses the quasi-isomorphisms of complexes
$$
L_{\mathcal{C}/\mathcal{B}} =
\Omega_{\mathcal{R}_\bullet/\mathcal{B}}
\otimes_{\mathcal{R}_\bullet} \mathcal{C}
\longrightarrow
\Omega_{\overline{\mathcal{R}}_\bullet/\mathcal{B}}
\otimes_{\overline{\mathcal{R}}_\bullet} \mathcal{C}
\longleftarrow
\Omega_{\overline{\mathcal{Q}}_\bullet/\mathcal{B}}
\otimes_{\overline{\mathcal{Q}}_\bullet} \mathcal{C}
$$
All the constructions used above can first be done on the level
of presheaves and then sheafified. Hence to prove sequences are exact,
or that map are quasi-isomorphisms it suffices to prove the corresponding
statement for the ring maps
$\mathcal{A}(U) \to \mathcal{B}(U) \to \mathcal{C}(U)$
which are known. This finishes the proof in the case that
$\mathcal{B} \to \mathcal{C}$ is injective.

\medskip\noindent
In general, we reduce to the case where $\mathcal{B} \to \mathcal{C}$ is
injective by replacing $\mathcal{C}$ by $\mathcal{B} \times \mathcal{C}$ if
necessary. This is possible by the argument given in
Remark \ref{remark-triangle} by
Lemma \href{https://stacks.math.columbia.edu/tag/08SY}{lemma compute L product sheaves rings (Tag 08SY)}.
\end{proof}
\end{document}
